\subsection{G 分解为双陪集}

定理 1. 有 G = BWB。映射 \(w \mapsto C(w)\) 是从 W 到关于 B 的 G 的双陪集集合 B\textbackslash G/B 的一个双射。

显然，BWB 在映射 \(x \mapsto x^{-1}\) 下不变，而引理 1 表明它在乘法下不变。由于它包含 B 和 N，故等于 G。

还需证明 \(\mathrm{C}(w) \neq \mathrm{C}(w')\)，当 \(w \neq w'\) 时。为此，我们将对整数 \(q\) 作归纳，证明以下断言：

(Aq) 若 w 和 \(w'\) 是 W 中满足 \(l_{S}(w) \geqslant l_{S}(w') = q\) 的不同元素，则 \(C(w) \neq C(w')\)。

（关于 \(l_{\mathbb{S}}(w)\) 的定义，见 §1，第 1 号。）

当 \(q = 0\) 时此断言显然，因为此时 \(w' = 1\) 且 \(w \neq 1\)，所以 \(C(w') = B\) 而 \(C(w) \neq B\)。

假设 \(q \geqslant 1\)，且 \(w, w'\) 满足 \((\mathbf{A}_q)\) 的假设。存在 \(s \in S\)，使得 \(sw'\) 的长度为 \(q - 1\)。我们有

\[
l _ {S} (w) > l _ {S} \left(s w ^ {\prime}\right)\tag{6}
\]

因此 \(w \neq sw'\)。此外，\(sw \neq sw'\)；由第 1 节第 1 号的公式 (3)，有

\[
l _ {\mathrm{S}} (s w) \geqslant l _ {\mathrm{S}} (w) - 1 \geqslant l _ {\mathrm{S}} (s w ^ {\prime}) = q - 1.\tag{7}
\]

由归纳假设，\(C(sw')\) 与 \(C(w)\) 及 \(C(sw)\) 均不相同；由公式 (2) 得

\[
\mathrm{C} (s w ^ {\prime}) \cap \mathrm{C} (s). \mathrm{C} (w) = \varnothing .\tag{8}
\]

由于 \(\mathrm{C}(sw') \subset \mathrm{C}(s).\mathrm{C}(w')\)，最后得到 \(\mathrm{C}(w) \neq \mathrm{C}(w')\)。

备注。上述证明没有使用公理 (T4)。

